In \cite{Beu13}, Beukers shows how to find a subgroup of the full monodromy group using Mellin--Barnes integral solutions of the associated $A$-hypergeometric system. This method only works under very restrictive conditions. These conditions are necessary to ensure the existence of a~basis of solutions in terms of Mellin--Barnes integrals. The monodromy groups found by Beukers' method are with respect to this Mellin--Barnes basis. In Sections~\ref{chap-ahyp}--\ref{chap-mon}, we will fix notation and introduce $A$-hypergeometric functions and Beukers' method.

The solution space of an $A$-hypergeometric system can be written in terms of integrals over generalised Pochhammer cycles \cite{Beu08} or over twisted cycles \cite{Kit94}. The unique invariant Hermitian form is the intersection form on these cycles. Actually computing this intersection form is computationally hard. Similarly it is in general computationally hard to explicitly compute an invariant hermitian form directly from a given monodromy group. In \cite{Beu89} Beukers and Heckman construct a unique invariant Hermitian form for the generalized hypergeometric system satisfied by ${}_{n+1}F_{n}$ using properties of its monodromy group. The goal of this paper is to give an explicit construction of the invariant Hermitian form over the monodromy group as constructed by Beukers' method \cite{Beu13}. The construction of this Hermitian form is given in Theorem \ref{MB-Hermitian} and its proof covers Sections~\ref{chap-herm} and~\ref{chap-res}. This invariant Hermitian form gives us insight into the structure of monodromy groups for $A$-hypergeometric functions with a Mellin--Barnes basis. The techniques provided in this paper may be used to extend these results to $A$-hypergeometric functions in general.
